\section{Introduction}
\label{sec:introduction}

\subsection{Background and Research Question}

Multi-factor models have been around since Fama and French introduced their three-factor framework in 1993. The hard part is not building the factors. The hard part is deciding how much weight to give each one, especially when market conditions change.

Our question: Can IC-based weighting provide a practical, interpretable approach to factor timing for Chinese A-shares?

\subsection{Data Statement}

We obtained real CSI 300 market data via Baostock API: 94 constituent stocks, 1,212 trading days (2020-01-01 to 2024-12-31). All results are from walk-forward backtesting with realistic costs. No simulation.

\subsection{Problem Statement}

Factor ICs vary over time. Static weights calibrated over long horizons misallocate when conditions shift. The 2022 Chinese market correction is a case in point: momentum factors that worked in 2020-2021 stopped working.

Three issues with existing approaches. First, static weights ignore IC dynamics. Second, slow-adapting rules cannot respond to rapid regime changes. Third, regime-agnostic weighting fails to exploit conditional factor effectiveness.

\subsection{Research Contribution}

We propose an IC-based dynamic weighting methodology with three components: rolling IC estimation, ICIR-based weighting, and risk control (drawdown limits, volatility targeting).

Our contributions: real CSI 300 validation (94 stocks, 1,212 days), correct Newey-West HAC statistics with Schwert lag selection, transparent reporting of marginal significance (p = 0.085 two-tailed, 0.0425 one-tailed), and documented factor ICs (momentum ICIR = 1.49, 96\% positive days).

The strategy achieved 14.79\% annualized return with 15.31\% maximum drawdown, Sharpe ratio 0.71. The statistics are marginal. The economics are meaningful.

\subsection{Paper Organization}

Section~\ref{sec:related} reviews related work. Section~\ref{sec:methodology} presents the IC-based weighting framework. Section~\ref{sec:results} presents real data results. Section~\ref{sec:discussion} discusses implications and limitations. Section~\ref{sec:conclusion} concludes with recommendations.
